\section{Introduction}

\subsection{Motivation}
Social media is integral to society with over 4.8 billion users worldwide\cite{umaine_socialmedia_stats} as it provides accessible communication and entertainment. However, most platforms like Facebook 
are centralised, meaning that a single provider controls identity, data storage and moderation. Thus raising censor censorship concerns (un-ethical banning of users), data privacy concerns and biases in what users see (algorithmitic bias) \cite{Ruwe2023}.

These limitations motivate interest in decentralised alternatives, decentralised means control is not in the handle of a single organisation but is distributed across users or nodes \cite{Ruwe2023}.The foundational principles of decentralisation emphasise transparency, user control and resilience against single points of failure. Transparency means that all actions including moderated decisions can be independently verified, user control allows users to have control over who can access their data and resilience prevents the system depending on a vulnerable central server. These principles have been established across decentralised identity and distributed system papers. \textbf{Achieving these goals together remains challenging, as existing approaches make different trade-offs in control, resilience and infrastructure requirements} \cite{w3c_did_core}.


\subsection{Problem Definition}
\textbf{Existing alternatives address different aspects of decentralisation, but vary in how they support user-controlled identity, replication, discovery and social interaction.} For example, the popular Mastadon ActivityPub protocol enables federation across servers \cite{wesleyac_identity_decentralization}, but identity still remains server-bound. The reason for this is that users depend on administrators who control domain names and the moderation policies. \textbf{An earlier thesis also identified backend dependencies as a challenge in attempts to make ActivityPub fully serverless.}

Similarly, protocols like Secure Scuttlebutt (SSB), Nostr and Authenticated Transfer (AT) Protocol decentralise social interactions through different models, including peer-to-peer exchange, relays and Personal Data Servers (PDS), each with its own infrastructure and portability trade-offs. \textbf{Git-based prototypes demonstrate that Git can support repository-based data exchange, although they vary in their features and infrastructure requirements.} \textbf{The approaches reviewed do not all provide the same combination of structured social actions, portable identity, discovery and a consistent social vocabulary.}

\textbf{Thus, the core problem is the lack of a lightweight, interoperable, Git-native protocol for social interactions that supports user-controlled identity and efficient replication without requiring a central platform to store and distribute social data.}


\subsection{Research Question}
\textbf{How can we design a lightweight, Git-native social protocol that provides portable identity, structured social actions, efficient replication and a complete social vocabulary, while supporting repository discovery without relying on a central social platform for social-data exchange?}


\subsection{Proposed Direction}
This project will conduct a systematic analysis of existing decentralised social media protocols, examining their architectures, strengths, limitations and relevance to the problem. 
This includes both non-Git based and Git-based attempts, providing concrete functional and non-functional requirements for a decentralised social protocol. 

\textbf{The design phase will reuse and extend strengths identified across both Git-based and non-Git-based protocols, while prioritising approaches that avoid dependence on centralised infrastructure for core social-data exchange.} Building upon these insights, the project will implement a Git-native prototype leveraging Git’s existing infrastructure to define identity models, replication flows and social vocabulary.

\subsection{Scope}
The project focuses on protocol design rather than building a full social network application. \textbf{The protocol aims to remain minimal and extensible, avoiding dependence on a central platform for social-data exchange while allowing supporting services for functions such as discovery.} The scope includes decentralised identity, efficient replication, repository discovery and a structured social vocabulary. Sophisticated User interface, moderation tooling and large-scale deployment are out of scope.
